% --- Archon physics formalization source begin ---
% archon:physics
% archon:covers PhyXMiniProblems/problem_phyx_mini_0524.lean
% archon:source-report reports/phyx_mini/problem_phyx_mini_0524.source.json

\chapter{Physics problem phyx\_mini\_0524}
\label{ch:PhyXMiniProblems_problem_phyx_mini_0524}

\paragraph{Problem source.}
\#\# Physical scenario

The speed of the Earth in its orbit is \textbackslash{}(29.8 \textbackslash{}, \textbackslash{}text\{km/s\}\textbackslash{}). If that is the magnitude of the velocity \textbackslash{}( \textbackslash{}vec\{v\} \textbackslash{}) of the ether wind in figure.

\#\# Question

Find the angle \textbackslash{}( \textbackslash{}phi \textbackslash{}) between the velocity of light \textbackslash{}( \textbackslash{}vec\{c\} \textbackslash{}) in vacuum and the resultant velocity of light if there were an ether.

\#\# Answer choices

- A: A: \textbackslash{}(6.95 \textbackslash{}times 10\textasciicircum{}\{-5\} \textbackslash{}

- B: B: \textbackslash{}(9.81 \textbackslash{}times 10\textasciicircum{}\{-5\} \textbackslash{}

- C: C: \textbackslash{}(6.32 \textbackslash{}times 10\textasciicircum{}\{-5\} \textbackslash{}

- D: D: \textbackslash{}(9.94 \textbackslash{}times 10\textasciicircum{}\{-5\} \textbackslash{}

\#\# Figure caption (auxiliary; use the image as primary evidence)

The image depicts a vector diagram with three vectors represented as red arrows. Here's a detailed breakdown:

1. **Vectors**:
   - The rightward horizontal vector is labeled \textbackslash{}(\textbackslash{}vec\{v\}\textbackslash{}).
   - The diagonal vector pointing upwards to the right is labeled \textbackslash{}(\textbackslash{}vec\{c\}\textbackslash{}).
   - A vertical vector completes the right triangle formed by \textbackslash{}(\textbackslash{}vec\{v\}\textbackslash{}) and \textbackslash{}(\textbackslash{}vec\{c\}\textbackslash{}).

2. **Angle**:
   - The angle between \textbackslash{}(\textbackslash{}vec\{c\}\textbackslash{}) and the vertical vector is labeled as \textbackslash{}(\textbackslash{}phi\textbackslash{}).

3. **Magnitude Text**:
   - To the left of the diagram, there's text indicating "Magnitude: \textbackslash{}(\textbackslash{}sqrt\{c\textasciicircum{}2 - v\textasciicircum{}2\}\textbackslash{})".

Overall, the diagram illustrates the geometric relationship between the vectors, likely relating to concepts in physics or mathematics such as velocity, light speed, or vector components.

\#\# Dataset metadata

Relativity; Spatial Relation Reasoning, Physical Model Grounding Reasoning

\paragraph{Recorded answer/context.}
D: D: \textbackslash{}(9.94 \textbackslash{}times 10\textasciicircum{}\{-5\} \textbackslash{}

\paragraph{Figure/image path.}
/root/proposal\_for\_physic/science-mango/phyx\_mini\_run/phyx\_data/test\_image/524.png

\paragraph{Formalization target.}
create a compiling Lean file with sorry bodies at `PhyXMiniProblems/problem\_phyx\_mini\_0524.lean`. The Lean declarations must preserve the physical quantities, dimensions or dimensional roles, figure labels, governing-law hypotheses, and final relation expressed by this problem.
Use Mathlib/Physlib names found through LeanExplore where available. If a physics API is missing, introduce faithful local abstractions rather than scalar placeholder aliases.

\begin{theorem}[Physics formalization target]
\label{thm:physics:phyx_mini_0524:target}
\lean{PhyXMiniProblems.ProblemPhyXMini0524.problem_phyx_mini_0524}
\uses{def:physics:phyx-mini-0524:phyxminiproblems-problemphyxmini0524-speedinkilometerspersecond, def:physics:phyx-mini-0524:phyxminiproblems-problemphyxmini0524-etherwindvelocitydiagram, def:physics:phyx-mini-0524:phyxminiproblems-problemphyxmini0524-matchesproblemspeeddata, def:physics:phyx-mini-0524:phyxminiproblems-problemphyxmini0524-hasphysicalvelocitymagnitudes, def:physics:phyx-mini-0524:phyxminiproblems-problemphyxmini0524-matchessuppliedetherwindfigure, def:physics:phyx-mini-0524:phyxminiproblems-problemphyxmini0524-satisfiesclassicalethervelocityaddition, def:physics:phyx-mini-0524:phyxminiproblems-problemphyxmini0524-answerchoice, def:physics:phyx-mini-0524:phyxminiproblems-problemphyxmini0524-displayedangleradians, def:physics:phyx-mini-0524:phyxminiproblems-problemphyxmini0524-roundstodisplayedangle}
The assigned autoformalize agent should translate this physics problem into Lean declarations in the covered file, with theorem and lemma proof bodies written as `by sorry`.
\end{theorem}
\begin{proof}
This is an autoformalization task, not a proof task. Produce statements that can later be proved by the physics prover without weakening the physical model.
\end{proof}

% --- Archon declaration topology begin ---
\section{Declaration topology}
\begin{definition}[velocity Dimension]
  \label{def:physics:phyx-mini-0524:phyxminiproblems-problemphyxmini0524-velocitydimension}
  \lean{PhyXMiniProblems.ProblemPhyXMini0524.velocityDimension}
  Velocity has physical dimension length divided by time
\end{definition}
\begin{proof}
  This is the declared model definition of velocity Dimension.
\end{proof}
\begin{definition}[Planar Velocity Quantity]
  \label{def:physics:phyx-mini-0524:phyxminiproblems-problemphyxmini0524-planarvelocityquantity}
  \lean{PhyXMiniProblems.ProblemPhyXMini0524.PlanarVelocityQuantity}
  \uses{def:physics:phyx-mini-0524:phyxminiproblems-problemphyxmini0524-velocitydimension}
  A unit-independent two-dimensional physical velocity vector
\end{definition}
\begin{proof}
  This is the declared model definition of Planar Velocity Quantity.
\end{proof}
\begin{definition}[Speed Quantity]
  \label{def:physics:phyx-mini-0524:phyxminiproblems-problemphyxmini0524-speedquantity}
  \lean{PhyXMiniProblems.ProblemPhyXMini0524.SpeedQuantity}
  \uses{def:physics:phyx-mini-0524:phyxminiproblems-problemphyxmini0524-velocitydimension}
  A unit-independent signed speed readout with velocity dimension
\end{definition}
\begin{proof}
  This is the declared model definition of Speed Quantity.
\end{proof}
\begin{definition}[planar Velocity Readout]
  \label{def:physics:phyx-mini-0524:phyxminiproblems-problemphyxmini0524-planarvelocityreadout}
  \lean{PhyXMiniProblems.ProblemPhyXMini0524.planarVelocityReadout}
  \uses{def:physics:phyx-mini-0524:phyxminiproblems-problemphyxmini0524-planarvelocityquantity}
  Read a planar velocity in coherent selected length and time units
\end{definition}
\begin{proof}
  This is the declared model definition of planar Velocity Readout.
\end{proof}
\begin{definition}[speed Readout]
  \label{def:physics:phyx-mini-0524:phyxminiproblems-problemphyxmini0524-speedreadout}
  \lean{PhyXMiniProblems.ProblemPhyXMini0524.speedReadout}
  \uses{def:physics:phyx-mini-0524:phyxminiproblems-problemphyxmini0524-speedquantity}
  Read a speed in coherent selected length and time units
\end{definition}
\begin{proof}
  This is the declared model definition of speed Readout.
\end{proof}
\begin{definition}[planar Velocity In Kilometers Per Second]
  \label{def:physics:phyx-mini-0524:phyxminiproblems-problemphyxmini0524-planarvelocityinkilometerspersecond}
  \lean{PhyXMiniProblems.ProblemPhyXMini0524.planarVelocityInKilometersPerSecond}
  \uses{def:physics:phyx-mini-0524:phyxminiproblems-problemphyxmini0524-planarvelocityquantity, def:physics:phyx-mini-0524:phyxminiproblems-problemphyxmini0524-planarvelocityreadout}
  Planar-velocity readout in the kilometres-per-second units of the source
\end{definition}
\begin{proof}
  This is the declared model definition of planar Velocity In Kilometers Per Second.
\end{proof}
\begin{definition}[speed In Kilometers Per Second]
  \label{def:physics:phyx-mini-0524:phyxminiproblems-problemphyxmini0524-speedinkilometerspersecond}
  \lean{PhyXMiniProblems.ProblemPhyXMini0524.speedInKilometersPerSecond}
  \uses{def:physics:phyx-mini-0524:phyxminiproblems-problemphyxmini0524-speedquantity, def:physics:phyx-mini-0524:phyxminiproblems-problemphyxmini0524-speedreadout}
  Speed readout in the kilometres-per-second units of the source
\end{definition}
\begin{proof}
  This is the declared model definition of speed In Kilometers Per Second.
\end{proof}
\begin{definition}[Velocity Arrow Label]
  \label{def:physics:phyx-mini-0524:phyxminiproblems-problemphyxmini0524-velocityarrowlabel}
  \lean{PhyXMiniProblems.ProblemPhyXMini0524.VelocityArrowLabel}
  The three arrows in the supplied velocity diagram
\end{definition}
\begin{proof}
  This is the declared model definition of Velocity Arrow Label.
\end{proof}
\begin{definition}[Ether Wind Velocity Diagram]
  \label{def:physics:phyx-mini-0524:phyxminiproblems-problemphyxmini0524-etherwindvelocitydiagram}
  \lean{PhyXMiniProblems.ProblemPhyXMini0524.EtherWindVelocityDiagram}
  \uses{def:physics:phyx-mini-0524:phyxminiproblems-problemphyxmini0524-planarvelocityquantity, def:physics:phyx-mini-0524:phyxminiproblems-problemphyxmini0524-speedquantity, def:physics:phyx-mini-0524:phyxminiproblems-problemphyxmini0524-velocityarrowlabel}
  The physical data represented in the image. Speed magnitudes are independent dimensionful quantities and are related to the velocity arrows only by the explicit physical assumptions below. In particular, anglePhiRadians is not assigned an answer-choice value here
\end{definition}
\begin{proof}
  This is the declared model definition of Ether Wind Velocity Diagram.
\end{proof}
\begin{definition}[Matches Problem Speed Data]
  \label{def:physics:phyx-mini-0524:phyxminiproblems-problemphyxmini0524-matchesproblemspeeddata}
  \lean{PhyXMiniProblems.ProblemPhyXMini0524.MatchesProblemSpeedData}
  \uses{def:physics:phyx-mini-0524:phyxminiproblems-problemphyxmini0524-speedinkilometerspersecond, def:physics:phyx-mini-0524:phyxminiproblems-problemphyxmini0524-etherwindvelocitydiagram}
  The numerical speed data stated or standard in the problem. The Earth/ether speed is 29.8 km/s = 149/5 km/s; the vacuum light-speed magnitude is Physlib's exact dimensionful constant
\end{definition}
\begin{proof}
  This is the declared model definition of Matches Problem Speed Data.
\end{proof}
\begin{definition}[Has Physical Velocity Magnitudes]
  \label{def:physics:phyx-mini-0524:phyxminiproblems-problemphyxmini0524-hasphysicalvelocitymagnitudes}
  \lean{PhyXMiniProblems.ProblemPhyXMini0524.HasPhysicalVelocityMagnitudes}
  \uses{def:physics:phyx-mini-0524:phyxminiproblems-problemphyxmini0524-planarvelocityreadout, def:physics:phyx-mini-0524:phyxminiproblems-problemphyxmini0524-speedreadout, def:physics:phyx-mini-0524:phyxminiproblems-problemphyxmini0524-planarvelocityinkilometerspersecond, def:physics:phyx-mini-0524:phyxminiproblems-problemphyxmini0524-speedinkilometerspersecond, def:physics:phyx-mini-0524:phyxminiproblems-problemphyxmini0524-velocityarrowlabel, def:physics:phyx-mini-0524:phyxminiproblems-problemphyxmini0524-etherwindvelocitydiagram}
  A speed magnitude is the Euclidean norm of its corresponding velocity-vector readout in every coherent unit system. Positivity and nonzero arrows record the nondegenerate physical branch of the diagram
\end{definition}
\begin{proof}
  This is the declared model definition of Has Physical Velocity Magnitudes.
\end{proof}
\begin{definition}[Matches Supplied Ether Wind Figure]
  \label{def:physics:phyx-mini-0524:phyxminiproblems-problemphyxmini0524-matchessuppliedetherwindfigure}
  \lean{PhyXMiniProblems.ProblemPhyXMini0524.MatchesSuppliedEtherWindFigure}
  \uses{def:physics:phyx-mini-0524:phyxminiproblems-problemphyxmini0524-speedreadout, def:physics:phyx-mini-0524:phyxminiproblems-problemphyxmini0524-planarvelocityinkilometerspersecond, def:physics:phyx-mini-0524:phyxminiproblems-problemphyxmini0524-etherwindvelocitydiagram}
  Literal and geometric information read from the primary raster. Coordinate 0 is horizontal and coordinate 1 is vertical. Thus v points left, the resultant points upward, and c points diagonally upward and right. The symbolic magnitude written beside the resultant is imposed in every coherent unit choice. None of these fields assigns the requested numerical angle
\end{definition}
\begin{proof}
  This is the declared model definition of Matches Supplied Ether Wind Figure.
\end{proof}
\begin{definition}[Satisfies Classical Ether Velocity Addition]
  \label{def:physics:phyx-mini-0524:phyxminiproblems-problemphyxmini0524-satisfiesclassicalethervelocityaddition}
  \lean{PhyXMiniProblems.ProblemPhyXMini0524.SatisfiesClassicalEtherVelocityAddition}
  \uses{def:physics:phyx-mini-0524:phyxminiproblems-problemphyxmini0524-planarvelocityreadout, def:physics:phyx-mini-0524:phyxminiproblems-problemphyxmini0524-etherwindvelocitydiagram}
  The governing classical ether model: velocities add by the parallelogram law. The orientation agrees with the image, where following c and then the leftward v reaches the head of the vertical resultant. This law is stated for every coherent unit choice and contains no angle value
\end{definition}
\begin{proof}
  This is the declared model definition of Satisfies Classical Ether Velocity Addition.
\end{proof}
\begin{definition}[Answer Choice]
  \label{def:physics:phyx-mini-0524:phyxminiproblems-problemphyxmini0524-answerchoice}
  \lean{PhyXMiniProblems.ProblemPhyXMini0524.AnswerChoice}
  Labels of the four dimensionless radian values printed in the source
\end{definition}
\begin{proof}
  This is the declared model definition of Answer Choice.
\end{proof}
\begin{definition}[displayed Angle Radians]
  \label{def:physics:phyx-mini-0524:phyxminiproblems-problemphyxmini0524-displayedangleradians}
  \lean{PhyXMiniProblems.ProblemPhyXMini0524.displayedAngleRadians}
  \uses{def:physics:phyx-mini-0524:phyxminiproblems-problemphyxmini0524-answerchoice}
  Radian readout printed beside each answer label
\end{definition}
\begin{proof}
  This is the declared model definition of displayed Angle Radians.
\end{proof}
\begin{definition}[recorded Dataset Answer]
  \label{def:physics:phyx-mini-0524:phyxminiproblems-problemphyxmini0524-recordeddatasetanswer}
  \lean{PhyXMiniProblems.ProblemPhyXMini0524.recordedDatasetAnswer}
  \uses{def:physics:phyx-mini-0524:phyxminiproblems-problemphyxmini0524-answerchoice}
  The source dataset records answer label D
\end{definition}
\begin{proof}
  This is the declared model definition of recorded Dataset Answer.
\end{proof}
\begin{definition}[Rounds To Displayed Angle]
  \label{def:physics:phyx-mini-0524:phyxminiproblems-problemphyxmini0524-roundstodisplayedangle}
  \lean{PhyXMiniProblems.ProblemPhyXMini0524.RoundsToDisplayedAngle}
  Agreement with a value printed to the displayed precision. Adjacent decimal places at this precision differ by 10⁻⁷, so the half-unit rounding radius is 5 * 10⁻⁸ = 1/20000000 radians
\end{definition}
\begin{proof}
  This is the declared model definition of Rounds To Displayed Angle.
\end{proof}
\begin{lemma}[ether Wind Angle Sine Relation]
  \label{lem:physics:phyx-mini-0524:phyxminiproblems-problemphyxmini0524-etherwindanglesinerelation}
  \lean{PhyXMiniProblems.ProblemPhyXMini0524.etherWindAngleSineRelation}
  \uses{def:physics:phyx-mini-0524:phyxminiproblems-problemphyxmini0524-speedinkilometerspersecond, def:physics:phyx-mini-0524:phyxminiproblems-problemphyxmini0524-etherwindvelocitydiagram, def:physics:phyx-mini-0524:phyxminiproblems-problemphyxmini0524-hasphysicalvelocitymagnitudes, def:physics:phyx-mini-0524:phyxminiproblems-problemphyxmini0524-matchessuppliedetherwindfigure, def:physics:phyx-mini-0524:phyxminiproblems-problemphyxmini0524-satisfiesclassicalethervelocityaddition}
  The right-triangle geometry first gives sin phi = v/c. This intermediate statement is derived from the norm, figure, and velocity-addition assumptions; it is not included in any premise structure
\end{lemma}
\begin{proof}
  The conclusion follows from the governing relations and figure readouts named in the statement of ether Wind Angle Sine Relation.
\end{proof}
% --- Archon declaration topology end ---
% --- Archon physics formalization source end ---
